\chapter{常用数学关系与推导约定}
\label{app:mathematics}

本附录汇集正文使用的概率运算、微分条件与最优性关系。各式中的随机性可以来自抽样、处理分配、训练数据或模拟情景。计算期望与方差时，应先确定随机变量及其概率分布；研究设计决定了哪些量固定、哪些量随重复观测而变化。

\section{期望、条件化与全概率}
对有限取值随机变量 \(Z\)，\(\E[Z]=\sum_z z\Prob(Z=z)\)。对可积的连续变量，
将求和改为关于其分布的积分。只要各期望存在，就有线性关系
\(\E[aX+bY]=a\E[X]+b\E[Y]\)。这一性质对相关变量同样成立。若 \(X,Y\) 独立且可积，另有 \(\E[XY]=\E[X]\E[Y]\)。

若事件 \(B_1,\ldots,B_k\) 构成样本空间的一个划分，且相关条件概率有定义，则
\[
\Prob(A)=\sum_{j=1}^k\Prob(A\mid B_j)\Prob(B_j).
\]
当 \(\Prob(A)>0\) 时，贝叶斯公式给出
\[
\Prob(B_j\mid A)=
\frac{\Prob(A\mid B_j)\Prob(B_j)}{\sum_l\Prob(A\mid B_l)\Prob(B_l)}.
\]
\subjectindex{beiyesigongshi}{贝叶斯公式}
第~\ref{sec:information-value}~节的检测算例正是先计算分母的信号概率，再计算
条件状态概率。敏感度是 \(\Prob(+\mid\theta=1)\)，阳性后的风险是
\(\Prob(\theta=1\mid+)\)，两者的条件事件不同，分别描述检测性能和检测后的风险。

全期望公式 \(\E[Z]=\E[\E(Z\mid X)]\) 说明可以先在每个条件层内计算均值，
再用目标总体的层权重汇总。第~\ref{sec:causal-standardization}~节的标准化即利用
这一运算结构。在一致性、条件可交换性和正值性成立时，观测条件均值可以识别对应的潜在结果均值。

\section{方差、协方差与平均的精度}
\subjectindex{xiefangcha}{协方差}

若二阶矩有限，方差为 \(\operatorname{Var}(X)=\E[(X-\E X)^2]\)，协方差为 \(\operatorname{Cov}(X,Y)=\E[(X-\E X)(Y-\E Y)]\)。将线性组合减去其期望，展开平方并逐项取期望，得到
\[
\operatorname{Var}\!\left(\sum_i a_iX_i\right)
=\sum_i a_i^2\operatorname{Var}(X_i)
+2\sum_{i<j}a_ia_j\operatorname{Cov}(X_i,X_j).
\]

独立同分布观测的协方差为零，因而样本均值的方差随样本量增加而下降。若单次观测的方差为 \(\sigma^2\)，则 \(\operatorname{Var}(\bar X)=\sigma^2/n\)。同一设备、客户或班组的记录通常具有组内相关性，应在方差计算中保留相应的协方差项。正文的集群设计效应和集成模型方差分解均由上述展开式得到。

条件方差分解为
\[
\operatorname{Var}(Y)=\E[\operatorname{Var}(Y\mid X)]
+\operatorname{Var}(\E[Y\mid X]).
\]
第一项描述层内波动，第二项描述层间均值差异。该分解描述分层、协变量调整及产品组合对总体波动的影响。这里的层间差异是条件均值的变异，其因果解释取决于研究设计。

\section{导数、约束与最优性}
可微函数在其定义域的无约束内点取得局部极值时，一阶导数必为零，但一阶条件本身
既不保证极小值，也不处理边界。对区间上的二阶可微实函数，若二阶导数在整个区间
严格为正，则函数严格凸，存在的驻点是该区间内唯一的全局极小值。只在某点二阶
导数为正，仅能结合该点的一阶条件判断严格局部极小值。
第~\ref{sec:queue-cost}~节先限制 \(\mu>\lambda\)，再验证二阶导数为正，因而能够
得到对应模型的唯一最优服务率。

在有约束问题中，应先检查可行性，再谈目标值。对最小化问题，任何可行方案给出
最优目标的上界；对最大化问题则给出下界。对偶可提供另一方向的界，二者相等时
构成最优性证明。数值结果应同时报告目标函数值、约束残差、求解状态及最优性界。

\section{量纲、权重与边界情形}
相加项必须具有可比较的单位。若损失以元计、时间以小时计，就需要明确的时间价值
系数才能相加。概率、比例和利用率没有物理量纲，但它们的分母所代表的对象仍需
一致。订单良率、工序通过率和设备效率分别对应不同对象，其汇总权重应由各自的分子、分母及计算单位确定。

边界情形是检验推导的有效方法。信息完全无用时，信息价值应为零；行动效果为零而
成本为正时，净收益应为负；目标函数可分离且不存在耦合约束时，优化可分解为逐对象选择；在正文的稳定 M/M/1 排队模型中，到达率从下方趋近服务率时，平均等待时间趋于无穷。这些结果可用于检查公式、参数和程序输出。

\chapter{计算复现与配套材料}
\label{app:reproduction}

配套材料包括案例数据、计算程序和自动化测试。案例数据用于重现正文的计算过程，程序按照给定的变量定义与参数运行，测试核查关键数值及边界情形。本附录统一说明运行方式、数据约定和结果核查方法。

\section{运行入口与版本记录}
配套仓库以项目根目录为工作目录，使用 Python 3.11 或以上版本。依赖按仓库的
环境说明安装。一个典型调用方式为
\begin{lstlisting}[language=bash,caption={从项目根目录运行配套案例}]
PYTHONPATH=src python -m bdm_decision.cases.week06_ab_testing
\end{lstlisting}
全书的数据与程序路径均相对于项目根目录。命令中的 \texttt{weekNN} 对应相应章节编号，输入数据位于 \path{data/sample/}，案例程序位于 \path{src/bdm_decision/cases/}。每次正式复现应记录代码版本、
数据文件哈希、依赖版本、随机种子、参数和输出日期。涉及真实企业数据时，还要记录
脱敏、授权范围和保留期限。

\begin{longtable}{P{0.09\linewidth}P{0.50\linewidth}P{0.29\linewidth}}
\caption{各章配套模块与实现范围；模块均位于 \texttt{bdm\_decision.cases}}\\
\toprule
\rowcolor{BDMTableHead}
章 & 模块名 & 基础实现\\
\midrule
\endfirsthead
\toprule
\rowcolor{BDMTableHead}
章 & 模块名 & 基础实现\\
\midrule
\endhead
1 & \path{week01_problem_framing} & 经营症状与问题定义。\\
2 & \path{week02_enterprise_data_model} & 企业对象与关系示例。\\
3 & \path{week03_sql_kpi}、\path{week03_semiconductor_kpi} & SQL/KPI 与制造指标的简化计算。\\
4 & \path{week04_data_quality} & 剖析、质量规则与失败记录。\\
5 & \path{week05_statistics_dashboard} & 描述统计与看板数据。\\
6 & \path{week06_ab_testing} & 均值差异、区间与实验报告。\\
7 & \path{week07_prediction} & 流失预测与基础评价。\\
8 & \path{week08_causality} & 分层、IPW、平衡与两期DID。\\
9 & \path{week09_optimization} & 小规模库存枚举优化。\\
10 & \path{week10_process_mining}、\path{week10_ai_infra_capacity} & 日志、变体、直接跟随关系与模型服务容量。\\
\bottomrule
\end{longtable}

各模块的计算范围见表。其中，制造指标程序使用等间隔每日记录计算平均 WIP，以记录数构造滚动窗口，并按产出件数汇总质量项。这些约定分别要求记录间隔一致、每日记录完整以及产品计量单位可比。对不等间隔状态记录，平均 WIP 应改用时间加权；对跨产品的设备效率，应按第~\ref{ch:week03-sql-kpi}~章定义的理想工时计算。

\section{从手算到程序的核查顺序}
复现首先应使用可以穷举或手算的小例子。核对对象数、分子分母、单位、排序和时间
范围之后，再扩大数据量。对 SQL，重点检查连接基数和聚合状态；对实验，重点检查
分配单元和标准误；对优化，重点检查约束及基准方案；对智能体，重点检查最终业务
状态与失败恢复，并记录各步骤的返回结果。

测试包括常规样本和边界样本。前者检查预期的计算结果，后者覆盖零分母、空集合、重复键、迟到事件、不可行约束、越权请求及响应丢失。结果状态应区分有效数值、未定义、授权拒绝和执行失败。例如，比例的分母为零时，该比例在数学上未定义；授权拒绝则表明请求未满足权限条件。

制造指标示例程序对零分母返回数值 \(0\)，这是该程序的输出约定。分析报告应另行检查分母，将此类结果标记为未定义，并从比例比较中排除。对缺失均值范围、连续产能成本、标准化效应和加急组合选择，仓库测试同时保留手算结果与程序结果，便于检查数值转录和计算实现。

\chapter{实证分析的报告方法}
\label{app:evidence}

\section{研究事实、机制解释与决策建议}

案例分析应明确区分来源材料中的事实、分析者提出的机制解释，以及在给定条件下形成的建议。企业公告可以用于确认产品部署、合作时间与实施范围；因果收益的估计则依赖处理组与对照组的可比性、结果定义和相应的识别假设。方法论文提供可引用的理论依据，企业研究还需说明这些条件在具体数据和制度安排中的含义。

例如，企业披露采用某平台组织数据目录，是关于实施情况的事实；目录通过降低数据发现与理解成本改善分析效率，是可进一步检验的机制；是否采购该平台，则取决于需求、成本、替代方案、安全要求和实施能力。第~\ref{ch:week02-enterprise-data-model}~章的 Atlan 案例介绍实施内容，第~\ref{ch:week08-causality}~章讨论平台成效的研究设计。

\section{比较对象与评价指标}

生产率研究应规定任务内容、成果质量及时间和成本的计量范围。个人撰写初稿的时间、团队形成可用成果的时间和组织的总成本是不同的结果变量。组织层面的测量还应计入复核、返工、协调、等待和后续处置。时间节省比例以原用时为分母，单位时间产出增幅以原产出率为分母，两者应分别计算。研究人工智能工具时，应保存模型、提示、工具权限及资料版本，以识别各项改动及其影响。

因果报告应明确目标总体、行动定义、起点、结果窗口、识别假设、估计量和不确定性。
优化报告应明确目标函数、量纲、可行域、输入来源、求解状态、敏感性和执行责任。
系统报告应明确成功标准、违规标准、失败成本、人工介入和恢复结果。
研究报告可以在统一的问题定义下分别呈现统计效果、经济收益和系统运行结果。

\section{来源更新与引用规则}
书目中的访问日期是在线材料的核查日期，不代替出版日期。经典论文按正式发表信息著录；预印本保留其性质；官方产品
文档按所引用页面或版本著录。动态页面可能更新，重要复现应另行保留合法可用的
版本记录。引用标准与制度文件时，应注明文件名称、发布机构、版本及所依据的条款或公开说明。

正文采用统一编号引文，书末按作者及题名字顺组织书目；同一来源跨章使用同一编号。
章末文献导读指出基础与进阶阅读之间的关系。索引按中文拼音和英文字母混合排序，
将关键概念、方法、缩略语及企业案例作为检索入口。引用用于追溯论证依据，索引用于
回到本书中的概念位置，两者承担不同的阅读功能。

\chapter{企业人工智能系统的容量实验}
\label{app:ai-infra-lab}

模型服务的占用时间与部署成本共同约束容量和经济可行性，相关关系见
第~\ref{sec:ai-service-operations}~节与第~\ref{ch:week09-optimization}~章。以下参数均为
假定值，用于分析指定负载下的服务时间、容量与费用。模型预填充与生成阶段的区分参考《深入理解 AI Infra》
\cite{aiinfra2026}；独立验收和迭代调整的方法参考《生成式软件工程》\cite{gse2026}。

\section{请求负载与时间分解}
一项问数任务使用基础提示1000个 token，检索4段证据，每段平均350个 token，
输出300个 token。输入共为2400个 token。在指定的假设负载下，预填充
速度设为6000 token/秒，生成速度设为70 token/秒，固定耗时0.5秒，
查询和其他工具各耗时0.6秒和0.4秒。串行执行的无排队时间为
\[
s=0.5+\frac{2400}{6000}+\frac{300}{70}+0.6+0.4
 \approx 6.186\text{ 秒}.
\]
这是所设负载的均值近似，不包含排队和重试。对推理服务的容量计算还需明确
服务位置的含义。本例假设整个任务在查询和工具等待时也占用一个位置；
若模型在等待期释放推理位置，则应分别测量推理与外部工具的占用时间，
对各类资源分别设置容量约束。

\section{成本阈值与容量可行性}
云端输入每百万 token 设为0.75美元，输出为3美元，查询每次为0.0002美元。
因而每请求费用是
\[
c_c=\frac{2400(0.75)+300(3)}{10^6}+0.0002
=0.0029\text{ 美元}.
\]
本地服务的固定成本设为每小时3美元，可变成本每请求0.0002美元。
两者的费用交点为
\(q^*=3/(0.0029-0.0002)\approx1111\) 请求/小时。设本地有四个等效服务位置，
平均占用上限为\(0.7\)，则按本例的服务时间得到设计容量
\(3600\times4\times0.7/6.186\approx1629\) 请求/小时。
费用阈值与容量上限共同定义候选区间；对峰值到达和服务时间波动，
需要测量或仿真请求完成时间的分布。

\section{程序复算与验收}
从项目根目录执行以下命令，即可获得所有假设与测算结果。
\begin{lstlisting}[language=bash,caption={模型服务容量案例的执行命令}]
PYTHONPATH=src python -m bdm_decision.cases.week10_ai_infra_capacity
PYTHONPATH=src python -m bdm_decision.cases.week10_ai_infra_capacity \
    --requests-per-hour 1500 --retrieved-passages 10
\end{lstlisting}
程序源码在\path{src/bdm_decision/cases/week10_ai_infra_capacity.py}。第一次运行应输出
2400个输入 token、约6.186秒占用时间、900请求/小时下的2.61与3.18美元
的小时费用。第二次运行将检索段数增加到10，输入变为4500个 token；
服务时间、云端费用和容量界限须按变化后的输入重新计算。

\begin{table}[htbp]
\centering\small
\caption{检索长度与请求量变化时的容量与费用估算}
\begin{tabular}{rrrrrr}
\toprule
\rowcolor{BDMTableHead}
检索段数 & 请求/小时 & 占用秒数 & 云端费用 & 本地费用 & 本地容量\\
\midrule
4 & 900 & 6.186 & 2.61 & 3.18 & 符合\\
4 & 1500 & 6.186 & 4.35 & 3.30 & 符合\\
10 & 1500 & 6.536 & 6.71 & 3.30 & 符合\\
10 & 1600 & 6.536 & 7.16 & 3.32 & 超出\\
\bottomrule
\end{tabular}
\end{table}
表中费用单位是美元/小时。检索段数从4增至10时，服务占用增至约6.536秒，
本地设计容量降至约1542请求/小时。1600请求/小时的费用虽然较低，
但已超出本例的平均占用上限。可以尝试增加服务位置，但固定成本、调度性能
与可行区间均需重新计算。

验收应分为数学核对和运行核对。前者用本附录的手算结果检查程序输出，
并对负的请求量和无效利用率执行失败测试。后者应在预定模型、并发配置
和负载下测量首输出时间、任务完成时间、日内请求峰值、失败与重试成本。
案例的简化容量公式不充当服务级别的正式证明。
